%! TeX root = 550.tex
\documentclass{article}
\usepackage{default_notes}

\begin{document}
  
\noindent \colorbox{green}{550.6}\\
\begin{align*}
  Arccos'(x) &= -\frac{1}{\sqrt{1-x^{2}}}\\ 
             &= -(1-x^{2})^{-\frac{1}{2}}\\ 
             &= - \sum_{n=0}^{+\infty}\frac{(-\frac{1}{2})(-\frac{3}{2})\dots(-\frac{1}{2}-n+1)}{n!}(-x^{2})^{n}\\ 
             &= - \sum_{n=0}^{+\infty}\frac{1\times 3 \times 5 \times \dots \times (2n-1)}{2^{n}n!}(-1)^{n}(-1)^{n}x^{2n}\\ 
             &= -\sum_{n=0}^{+\infty}\frac{(2n)!}{2^{2n}(n!)^{2}}x^{2n}\\ 
             &= -\sum_{n=0}^{+\infty}\binom{2n}{n}\frac{x^{2n}}{4^{n}}\\ 
  Arccos(x) &= \frac{\pi}{2} - \sum_{n=0}^{+\infty}\binom{2n}{n}\frac{x^{2n+1}}{(2n+1)4^{n}}
\end{align*}

\noindent \colorbox{green}{550.3}\\ 
\begin{align*}
    \ln(1+t) = \sum_{n=1}^{+\infty}\frac{(-1)^{n-1}}{n}t^{n}
\end{align*}
Théorème radial d'Abel + continuité du logarithme et c'est cool \\

\noindent \colorbox{green}{550.4}\\ 
\begin{align*}
  \int_{0}^{1}\frac{dt }{1+t} &= \int_{0}^{1}\sum_{n=0}^{+\infty}(-t)^{n} dt
\end{align*}
\((t \mapsto \sum_{n=0}^{N}(-t)^{n})_{N \in \mathbb{N}}\)
converge simplement sur \(]-1, 1[\) vers \(t \mapsto 0\)\\ 
\underline{Dominons}
\begin{align*}
  |\sum_{n=0}^{+\infty}(-t)^{n}| &= |\frac{1-(-t^{n+1})}{1+t}| \\
                                 &\leq \frac{1+t}{1+t} = 1 \indent \text{indépendant de N, intégrable sur} ]0, 1[
\end{align*}
\begin{align*}
  \int_{0}^{1}\sum_{n=0}^{+\infty}(-t)^{n}dt &\xrightarrow[N \to +\infty]{} \int_{0}^{1}\sum_{n=0}^{+\infty}(-t)^{n}dt \\ 
  \sum_{n=0}^{N}\int_{0}^{1}(-t)^{n}dt &= \sum_{n=0}^{N}\frac{(-1)^{n}}{n+1}\\ 
                                       &\xrightarrow[n \to +\infty]{} \sum_{n=0}^{+\infty}\frac{(-1)^{n}}{n+1}\\ 
                                       &= \sum_{n=1}^{+\infty}\frac{(-1)^{n-1}}{n}
\end{align*}

\noindent \colorbox{green}{550.21}\\ 
Pour \(n \in \mathbb{N}^{*}\)
\begin{align*}
    I_{n}= \int_{1}^{+\infty} e^{-t^{n}}dt \\ 
\end{align*}
* Soit \(t \in ]1, +\infty[\)
\begin{align*}
  e^{-t^{n}} \xrightarrow[n \to +\infty]{} 0 
\end{align*}
* Dominons, \(\forall n \in \mathbb{N}^{*}, \forall t \in ]1, +\infty[\)
\begin{align*}
  |e^{-t^{n}}| &\leq e^{-t}
\end{align*}
\(e^{-t}\) indépendant de n, intégrable sur \(]1, +\infty[\)\\
Par convergence dominée
\begin{align*}
  I_{n} \xrightarrow[n \to +\infty]{} \int_{1}^{+\infty} 0 dt = 0
\end{align*}
b/ 
\begin{align*}
  t &= u^{\frac{1}{n}}\\ 
  dt &= \frac{1}{n}u^{\frac{1}{n}-1}\\
  I_{n}&=\int_{1}^{+\infty} e^{-u }\frac{1}{n}u^{\frac{1}{n}-1}du \\
       &= \frac{1}{n}\underbrace{\int_{1}^{+\infty} \frac{e^{-u }}{u}u^{\frac{1}{n}}du}_{J_{n}}\\ 
  J_{n}&= \int_{1}^{+\infty}\frac{e^{-u }}{u}u^{\frac{1}{n}} du 
\end{align*}
Soit u \(\in [1, +\infty[\)\\ 
\begin{align*}
  u^{\frac{1}{n}} = e^{\frac{1}{n}\ln u} &\xrightarrow[n\to +\infty]{} 1 \\ 
  \frac{e^{-u }}{u}u^{\frac{1}{n}} &\xrightarrow[n \to +\infty]{} \frac{e^{-u }}{u}
\end{align*}
Dominons, \(n \in \mathbb{N}^{*}, u \in ]1, +\infty[\)
\begin{align*}
  |\frac{e^{-u }}{u}u^{\frac{1}{n}}| &\leq \frac{e^{-u }}{u} u \\ 
                                     &= e^{-u } \text{ indépendant de n, intégrable sur } ]1, +\infty[
\end{align*}
\(I_{n}\) est une suite décroissante, positive, de limite nulle. 
Donc \(\sum_{}^{} I_{n}\) converge \\

\noindent \colorbox{green}{550.5}\\ 
a. 
\begin{align*}
  |\frac{1}{(2n+2)!}(2n)!| &= \frac{1}{(2n+2)(2n+1)}\\ 
                           &\xrightarrow[n \to +\infty]{} 0
\end{align*}
Donc d'après la règle d'Alembert \fbox{\(R = +\infty\)}\\ 
b. 
\begin{align*}
    ch x = \sum_{n=0}^{+\infty} \frac{x^{2n}}{(2n)!}
\end{align*}
c. 
\begin{align*}
    S(x) = \sum_{n=0}^{+\infty}\frac{x^{n}}{(2n)!}
\end{align*}
\underline{Si x = 0} S(0) = 1\\ 
\underline{Si \(x > 0\)} 
\begin{align*}
    x^{n} = (\sqrt{x})^{2n}\\ 
    S(x) = ch(\sqrt{x})
\end{align*}
\underline{Si \(x < 0\)}\\ 
\begin{align*}
  S(x) &= \sum_{n=0}^{+\infty}\frac{x^{n}}{(2n)!}\\ 
       &= \sum_{n=0}^{+\infty} \frac{(-1)^{n}(-x)^{n}}{(2n)!}\\ 
       &= \sum_{n=0}^{+\infty} \frac{(-1)^{n}(\sqrt{-x})^{2n}}{(2n)!}\\ 
       &= \cos(\sqrt{-x})
\end{align*}
\(\forall x, f(x) = S(x)\) qui est \(\mathcal{C}^{\infty}\) 
comme somme de série entière\\ 

\noindent \colorbox{green}{550.1}\\ 
\begin{align*}
  |\frac{1}{n(n+1)}| \leq \frac{1}{n^{2}} \text{ Donc converge }
\end{align*}
\begin{align*}
  S(x) &= \sum_{n=1}^{+\infty} \frac{x^{n}}{n} - \sum_{n=1}^{+\infty} \frac{x^{n}}{n+1} \\ 
       &= -\ln(1-x) - \frac{1}{x}(-\ln(1-x)-x)\\ 
       &= \ln(1-x)(\frac{1}{x}-1)+1\\ 
  S(x) &= 2\ln(\frac{2}{3}) + 1
\end{align*}

\noindent \colorbox{green}{550.2}\\ 
\begin{align*}
    \sum_{n=0}^{+\infty} \frac{(-1)^{n}}{2n+1}
\end{align*}
\((\frac{1}{2n+1})_{n \in \mathbb{N}}\) positive, décroissante,
de limite nulle donc la série converge \\ 
On pose \(S(x) = \sum_{n=0}^{+\infty}(-1)^{n} \frac{x^{2n+1}}{2n+1}\)
Donc x \(\in ]-1, 1[\)
\begin{align*}
  S(x) &= Arctan x + 0 \\ 
    S(x) &\xrightarrow[x \to 1]{} S(1) \text{ Par le théorème radial d'Abel} \\ 
         &= \frac{\pi}{4}
\end{align*}

\noindent \colorbox{green}{550.22}\\
c.
\begin{align*}
    \sum_{n=0}^{+\infty}(n^{2}+1)2^{n+1}x^{n}\\ 
    &= 2 \sum_{n=0}^{+\infty} (n+1)(n+2)(2x)^{n} - 3.2 \sum_{n=0}^{+\infty}(n+1)(2x)^{n} + 2.2 \sum_{n=0}^{+\infty}(2x)^{n} \\ 
    &= \frac{4}{(1-2x)^{3}}-\frac{6}{(1-2x)^{2}} + \frac{4}{1-2x}
\end{align*}
a. On note \(c_{n} = \sum_{k=1}^{n}\frac{1}{k} = \sum_{k=0}^{n}a_{k}b_{n-k}\)
avec \(
\begin{cases}
  a_{k} = \frac{1}{k} \text{ si k } \neq 0, k = 0 \text{sinon}
  b_{k} = 1
\end{cases}
\)
\begin{align*}
    S(x) = -\ln(1-x) \frac{1}{1-x}
\end{align*}

\noindent \colorbox{green}{550.27}\\
a. 
\begin{align*}
    x + \sum_{n=2}^{+\infty}(-1)^{n}(\frac{1}{n-1}-\frac{1}{n})x^{n}
\end{align*}
















\end{document}
